\SourceLeaf{244}%
1829---1830 finden lassen dürfte. Ich kann Ihnen aber den Titel einer anderen Schrift nennen, die Sie zweifellos sehr leicht über den Buchhandel erhalten können und die nur 4 Druckbogen umfasst.

\begin{center}
»Geometrische Untersuchungen zur Theorie der Parallel-Linien\\
von \textsc{Nicolaus} \textsc{Lobatschewsky} Kais. Russischem Staatsrath \emph{etc.}\\
Berlin 1840 in der G. Finckeschen Buchhandlung.«
\end{center}

Ich erinnere mich, damals in \textsc{Gersdorf}s Repertorium eine sehr abfällige Rezension dieses Buchs gelesen zu haben. Jeder etwas sachkundige Leser konnte der Rezension ansehen, dass sie von einem völlig unwissenden Menschen stammte. Seit ich Gelegenheit hatte, diese kleine Schrift selbst einzusehen, muss ich sie \textls{sehr günstig} beurteilen. Sie hat insbesondere viel mehr Eleganz und Präzision als die längeren Aufsätze von \textsc{Lobatschewski}. Diese gleichen eher einem verworrenen Wald, in dem es schwer ist, einen Weg und Überblick zu finden, ohne zuvor jeden Baum einzeln kennengelernt zu haben.

Über die in \textsc{Crelle} 17 S. 303 angeführte experimentelle Begrenzung habe ich aber in der Schrift von 1840 nichts gefunden. Ich werde mich daher wohl entschließen müssen, deswegen einmal an Herrn \textsc{Lobatschewski} selbst zu schreiben, dessen Aufnahme als korrespondierendes Mitglied unserer Sozietät ich vor einem Jahr veranlasst habe. Vielleicht schickt er mir dann den Kasaner Boten. Es wäre aber möglich, dass sich in den folgenden Jahrgängen der Kasaner Gelehrten Schriften von 1836---1838, die auch längere Aufsätze von \textsc{Lobatschewski} enthalten, etwas dazu findet. Ich besitze sie zwar, habe mich aber aus dem in meinem letzten Brief genannten Grund bisher nicht näher mit ihnen vertraut gemacht.

In seinem Dankesbrief für die Aufnahme in die Sozietät schrieb mir \textsc{Lobatschewski} damals übrigens, dass ihn seine vielen Verwaltungsaufgaben (er scheint Rector perpetuus der Universität zu sein) inzwischen ganz von der wissenschaftlichen Arbeit abgehalten hätten.

Für heute schließe ich mit der Bitte, mich bald wieder mit einigen Zeilen zu erfreuen.

\begin{flushright}
Ihren treu ergebenen\\
\textsc{C. F. Gauss}.
\end{flushright}

\par\bigskip
